Information Energetics: The Thermodynamics of Persistence at Cosmic Scale

The Reframed Narrative Arc

"How does a finite-capacity information-processing system maintain persistence as its internal complexity grows over cosmic time?"

I. The Universe as a Persistent System
Any system that persists must minimize Entropic Action (Papers 1–2)
The universe is the ultimate closed system: no external environment, no entropy sink, no outside observer
Therefore, its dynamics are governed by internal dissipation minimization, not external boundary conditions
This inverts the standard cosmological narrative: the universe is not "expanding into" anything; it is maintaining coherence against its own accumulated complexity
II. The Primordial State: Genesis as Protocol Establishment
Inflation is not a "field" but the handshake protocol establishing causal connectivity across the lattice
The spectral index nₛ = 1 − 1/(2ν) is the compression loss of this protocol — the information cost of synchronizing N=32 channels into ν=16 active ones
The null tensor mode r ≡ 0 is not a "prediction" but a structural necessity: no load, no gravitational radiation source (gravity is emergent load response, not primordial)
Baryogenesis is Reversible Encoding at cosmic scale: the Jarlskog mismatch (J ≈ 3×10⁻⁵) cannot be erased, so it is stored as persistent structural bias (matter vs. antimatter)
III. The Aging Universe: Systemic Latency
Early universe (radiation-dominated): low structural load, high "clock speed" — the unloaded rate H₀^early ≈ 67.4
Late universe (matter-dominated): high structural load, system slows to maintain coherence — the loaded rate H₀^late = H₀^early × L
The Lattice Load Factor L = 1 + S_Q + S_T is the thermodynamic overhead of maintaining causal order under load:
S_Q = 1/(2ν): digital overhead (finite bit-depth)
S_T = χ/∆: analog overhead (continuous boundary vs. discrete resonance)
The Hubble tension is not a discrepancy between measurements. It is the same system measured in two different operating regimes — idle vs. loaded.
IV. The Terminal State: Black Holes as Buffer Overflow
Black holes are not singularities. They are saturation events — local regions where information density K(r) → ν, causing channel capacity exhaustion
The event horizon is the Governor boundary: the system prevents runaway by freezing updates (n(r) → ∞, c' → 0)
The information paradox is resolved by recognizing that erasure is thermodynamically prohibited. The horizon is a write-protect surface; information is buffered, not destroyed
Hawking radiation is reflected information — the system's attempt to maintain impedance matching by releasing stored metadata
Bekenstein-Hawking entropy S = A/(D·ℓ_p²): the 1/4 factor is not arbitrary but the dimensionality cost of projecting 3D information onto a 2D surface (the holographic principle as capacity constraint)
V. The Dark Sector: Unallocated Capacity
Dark matter is not particles. It is geometric mass — the E8 mirror sector (248 − 48 = 200 dimensions) that cannot couple to the 4D causal manifold because it lacks temporal generators (no Lorentzian signature)
Dark energy is not a field. It is entropic exhaust — the minimum resolvable energy of the lattice (ρ_vac) modulated by load factor L(z)
The S₈ tension is damping from overloaded protocol: structure cannot grow when the coordination bandwidth is saturated


| Current Section                    | Revised Framing                                                 |
| ---------------------------------- | --------------------------------------------------------------- |
| "The Mechanism: Chiral Truncation" | "Genesis: Establishing the Protocol Layer"                      |
| "The Spectral Index"               | "Map Compression: The Lossy Encoding of Primordial State"       |
| "Tensor Modes"                     | "Governor Constraint: No Load, No Gravity"                      |
| "The Origin of Matter"             | "Reversible Encoding: The Jarlskog Mismatch as Persistent Bias" |
| "The Unloaded Universe"            | "Idle State: Baseline Clock Speed"                              |
| "The Loaded Universe"              | "Operating Under Load: Systemic Latency"                        |
| "The Hubble Tension"               | "Vital Sign Discrepancy: Idle vs. Loaded Readings"              |
| "Black Holes"                      | "Saturation Events: Buffer Overflow and Information Buffering"  |
| "Dark Sector"                      | "Unallocated Capacity: The Mirror Sector and Entropic Exhaust"  |


Conclusion:
IE cosmology: "The universe has these parameters because any system with finite capacity minimizing entropic action would exhibit them. The universe is not special; it is generic. Any sufficiently complex persistent system — biological, economic, computational, cosmic — must show: load-dependent latency, lossy compression of historical state, saturation events at capacity limits, and unallocated capacity appearing as 'dark' sectors."